% !TEX program = xelatex
% ChabotAI伙伴 V1.4.3 软件设计说明书（软件著作权登记·文档鉴别材料）
\documentclass[zihao=5,a4paper]{ctexart}
\usepackage[left=2.2cm,right=2.2cm,top=2.4cm,bottom=2.2cm,headheight=14pt,headsep=8pt]{geometry}
\usepackage{fancyhdr}
\usepackage{tabularx}
\usepackage{booktabs}
\usepackage{enumitem}

\pagestyle{fancy}
\fancyhf{}
\fancyhead[L]{\small ChabotAI伙伴 V1.4.3}
\fancyhead[R]{\small 软件设计说明书}
\fancyfoot[C]{\small 著作权人：陈聪\qquad 第 \thepage\ 页}
\renewcommand{\headrulewidth}{0.4pt}

\setcounter{tocdepth}{2}

\begin{document}

% ==================== 封面 ====================
\begin{titlepage}
  \centering
  \vspace*{1.6cm}
  {\Huge\bfseries ChabotAI伙伴}\\[0.6cm]
  {\Large AI 人格聊天应用软件}\\[1.2cm]
  {\huge\bfseries 软件设计说明书}\\[0.4cm]
  {\large V1.4.3}\\[1.6cm]

  \begin{tabular}{r@{\quad}l}
    \toprule
    软件全称 & ChabotAI伙伴 \\
    软件简称 & AI 伙伴 \\
    软件版本 & V1.4.3 \\
    著作权人 & 陈聪 \\
    文档类型 & 软件设计说明书 \\
    编制日期 & 2026 年 9 月 \\
    \bottomrule
  \end{tabular}

  \vspace{1.6cm}
  \begin{minipage}{0.86\textwidth}
    \zihao{-5}
    \textbf{【版权声明】}本软件设计说明书为“ChabotAI伙伴 V1.4.3”的配套文档，
    仅供计算机软件著作权登记审查使用。软件的著作权人为陈聪，受《中华人民共和国著作权法》
    及相关法律法规保护。未经著作权人书面许可，任何单位和个人不得以任何方式复制、
    传播、销售或用于其他商业用途。本说明书所描述的功能、结构、接口与数据设计
    均以实际交付的源程序为准。
  \end{minipage}

  \vspace{1.2cm}
  \begin{minipage}{0.86\textwidth}
    \zihao{-5}
    \textbf{【文档说明】}本文档依据软件源代码与功能实现编写，全面介绍软件的概述、
    运行环境、系统架构、模块功能设计、数据结构设计、接口设计、技术特点、测试结果
    与使用方法，供登记审查与后续维护参考。
  \end{minipage}

  \vfill
  \vspace*{0.8cm}
\end{titlepage}

% ==================== 目录 ====================
\begingroup
\zihao{-5}
\tableofcontents
\endgroup

\vspace{0.5em}

% ==================== 正文 ====================
\section{概述}

\subsection{软件简介}
ChabotAI伙伴是一款基于大语言模型的跨端 AI 人格聊天应用软件，无后端、纯端侧运行。
软件采用 uni-app（Vue 3）技术栈开发，一套源代码即可编译运行于 Android、iOS（App）、
H5 浏览器以及微信小程序等平台。软件调用任意 OpenAI 兼容的 /chat/completions 接口
完成智能对话，通过系统提示词定义 AI 伙伴的个性化人格，并围绕记忆、情景、上下文压缩、
表情包、语音阅读、拟真聊天等能力，为用户提供可持续、可定制的虚拟陪伴体验。

软件整体采用组件化与模块化设计：聊天页由若干子组件构成；底层逻辑按职责划分为
聊天服务（chat）、存储（storage）、记忆（memory）、提示词（prompts）、表情（emojis）、
语音（tts）、调试日志（log）、通知（notify）、导出（export）等多个功能模块。
模块之间通过统一门面与共享状态协同，职责清晰、易于扩展与维护。

\subsection{开发目的}
软件旨在为用户提供一款无需后端服务器、可自由接入任意大模型接口的 AI 人格聊天工具。
用户可自定义伙伴的人格、知识记忆与互动频率，使其逐渐成长为贴合用户习惯的个性化
AI 伙伴。软件同时强调端侧数据隐私：聊天记录、记忆与设置均保存在用户本地设备，
不上传任何云端。

\subsection{面向领域与行业}
软件面向人工智能与移动应用领域，适用于大模型对话、虚拟陪伴、数字人聊天、角色扮演
等场景，可供普通用户作为个人聊天工具使用，也可作为大模型端侧应用的参考实现。

\section{运行环境}

\subsection{硬件环境}
\textbf{开发硬件环境：}PC 兼容机一台，CPU 为 Intel Core i5 及以上，内存 16GB，
固态硬盘 512GB，配 Android 真机或模拟器用于调试验证。

\textbf{运行硬件环境：}Android 5.0 及以上 / iOS 9 及以上智能手机（内存 4GB 及以上），
或安装现代浏览器的 Windows / macOS 电脑。

\subsection{软件环境}
\textbf{开发软件环境：}操作系统为 Windows 10 64 位；开发工具为 HBuilderX、
Visual Studio Code；开发框架为 uni-app（Vue 3）；辅助工具链为 Node.js
（用于运行自动化测试脚本）。

\textbf{运行软件环境：}Android/iOS 原生运行环境、HTML5 浏览器内核、微信小程序
运行环境。软件本身无第三方运行时依赖，运行支撑软件主要为远程 OpenAI 兼容的
大模型接口服务与 Qwen-TTS 语音合成接口服务。

\section{系统架构}

\subsection{总体架构}
软件整体为单一应用工程，由四个层次构成：
\begin{enumerate}[itemsep=0pt,parsep=0pt]
  \item \textbf{页面视图层：}负责各页面与用户交互，包括聊天页、记忆管理页、
        表情管理页、设置页及若干设置子页；
  \item \textbf{聊天服务层：}封装发送、校验、记忆、情景、压缩、拟真等核心业务流程；
  \item \textbf{存储门面层：}统一封装跨端持久化读写；
  \item \textbf{外部接口层：}通过统一网络请求组件调用远程 LLM 接口与 TTS 语音合成接口。
\end{enumerate}

聊天页作为主界面持有会话状态，七个聊天子组件（头部、消息列表、输入栏、表情栏、
情景编辑、历史、人格）通过属性下发与事件上报机制与页面协调。

\subsection{模块划分}
程序主要模块划分如下：
\begin{itemize}[itemsep=0pt,parsep=0pt]
  \item \textbf{聊天服务域：}chat（门面与发送主链路）、chat-state（共享状态）、
        chat-settings（设置域）、chat-conversations（会话域）、chat-compress（压缩域）、
        chat-proactive（拟真聊天域）；
  \item \textbf{存储域：}storage（持久化门面）、storage-state（底层与共享状态）、
        storage-conversations（会话）、storage-scene（情景）、storage-api（API 预设）、
        storage-background（聊天背景）；
  \item \textbf{记忆域：}memory（记忆核心）、memory-similarity（相似度）、
        memory-parse（Memory 行解析）、memory-constants（常量）；
  \item \textbf{提示词域：}prompts（系统提示词构建、人格预设、接口预设）；
  \item \textbf{辅助能力域：}emojis（表情数据层）、llm（LLM 客户端）、tts（语音阅读）、
        log（调试日志）、notify（系统通知）、export（记录导出）；
  \item \textbf{页面层：}pages/chat 下聊天页与子组件、pages/memory 记忆页、
        pages/emoji 表情页、pages/settings 下设置页与设置子页。
\end{itemize}

\section{模块功能设计}

\subsection{智能聊天模块}
聊天模块负责对话主链路。发送消息时依次完成：检索记忆、组装系统提示词（人格、规则、
记忆指南、情景指南、当前状态、记忆上下文、记忆容量状态、表情清单）、注入未压缩历史、
调用 LLM、格式校验、情景与记忆更新、回复落库、执行维护与上下文压缩。

回复格式校验为聊天模块的关键设计：要求回复内含对话文本，Scene 行有内容，
Memory 行可解析，并校验表情占位均在表情清单内；不合格自动重新请求，
最多重试设置项指定的次数，从而保证落库数据规范。

若当前会话开启拟真聊天且时间模式为现实时间，则所有回复还需满足每条不超过一句、
连续表情不超过两个的约束，并按换行拆分为多条气泡落库。

\subsection{记忆系统模块}
记忆系统采用三级记忆模型：L1 核心事实永不衰减（重点事实永久保留）；L2 情景记忆
慢衰减，用作与用户的约定与待办清单；L3 临时信息快衰减，用于近期琐事与临时状态。

记忆由大模型驱动写入：系统提示词内置记忆书写指南，大模型在回复末尾输出结构化的
Memory 行，程序解析并执行新增、修改或删除操作；修改与删除按内容逐字匹配，
找不到匹配则静默忽略。

记忆具备如下核心机制：
\begin{itemize}[itemsep=0pt,parsep=0pt]
  \item \textbf{去重合并：}字符 bigram Jaccard 相似度与 LCS 序列相似度加权判定，
        超过阈值视为近似记忆并合并；
  \item \textbf{召回冷却：}记忆被召回后一段时间内禁止重复保存，复读视为强调并
        提升重要性（上限 5 级）；
  \item \textbf{分层检索：}全量召回 L1 核心事实、新鲜槽与 MMR 多样性槽混合检索，
        总配额 30 条、多样性系数 0.7；
  \item \textbf{周期维护：}L3 过期清理、L2/L3 升降级、容量淘汰（上限 200 条）；
  \item \textbf{时间模式：}支持现实时间与虚拟时间两种衰减模式，虚拟模式下按剧情
        时刻相对衰减，现实中断不衰减。
\end{itemize}

\subsection{会话管理模块}
软件支持多会话管理：可新建、切换、删除、复制会话；开始新对话时会复制当前会话的
完整设置快照；每个会话独立保存人格与设置；会话标题自动取首条用户消息；
提供历史弹窗进行切换、删除、压缩上文与导出对话记录；支持复制会话或仅复制记忆
到新会话。

\subsection{上下文压缩模块}
为控制请求长度、节省 token 并保持上下文连贯，软件支持将历史对话压缩为概要。
可手动触发，也可按设置条数在每轮回复落库后自动触发；压缩范围按上次进度计算，
保留最近指定条数的消息；超过批处理上限时分批调用、逐批合并概要，避免单次请求过大；
概要写入会话并注入后续请求，原始消息仍完整保留可翻阅。

\subsection{情景感知模块}
软件提供当前情景识别能力。大模型每次回复输出 Scene 行描述用户当前所处的场景，
程序结合注入的时间判断用户此刻在做什么，并持久化到情景历史（最多 10 条，先进先出）。
支持实时时间与虚拟时间两种模式：实时时间模式注入真实当前时间供情景判断；
虚拟时间模式不发送真实时间，由大模型自由想象情景氛围。聊天页顶部以情景条展示，
可点击查看、修改或清除。

\subsection{表情包系统模块}
软件内置自定义表情包系统：支持批量上传图片并逐张命名，命名全局唯一、限长；
支持改名、删除与拖拽排序；表情图片按平台持久化（App 与小程序压缩后存文件，
H5 压缩为 base64）。

对话中以 \$表情名\$ 占位引用表情并渲染为图片；程序提供消息拆分解析（跳过空白段、
裁剪首尾空白），清单外的未知占位原样保留；可通过设置开关控制是否将表情清单注入
系统提示词，以决定大模型是否主动使用表情。

\subsection{语音阅读模块}
软件对接 Qwen-TTS 非流式语音合成接口，可将 AI 回复朗读出来。支持音色选择
（内置常用音色清单并支持手动输入）、模型与接口地址配置、接口测试。

语音阅读具备多项细节设计：表情不朗读（仅朗读文本段）；文本截断到接口上限；
拿到音频直链后直接播放、不落盘；同一时刻只播一条；App 端等待可播放事件后再播放
并设兜底重试；首次播放失败会自动重新合成一次。

\subsection{拟真聊天模块}
拟真聊天用于模拟真人发消息：AI 会在设定的时间窗口内按频率档位随机主动给用户发
消息，每条不超过一句、连续表情不超过两个。

调度采用单条链式定时器，空闲零唤醒；退后台不清除定时器，App 真后台脚本挂起时由
回前台补发；到期触发时刻持久化到存储，进程重启后对齐恢复；触发具备门禁
（开关、现实时间、未抑制、时段窗口、已有用户消息），每次被拦截都会记录原因日志。
可选开启系统后台通知（App 本地通知 / H5 浏览器通知），仅在应用处于后台时弹出。

\subsection{设置与数据管理模块}
设置页集中管理：接口配置（接口地址、API Key、模型、温度，支持最多 3 套预设快速
切换）、思考模式、请求次数上限、人格配置、聊天表情包开关、语音阅读配置、聊天背景图、
上下文压缩间隔、数据管理（清空、迁移、恢复）与调试日志面板。

软件内置调试日志系统：采用环形缓冲区，记录请求、响应、错误与信息四类日志，
可在调试面板按类型过滤、展开查看详情、刷新与清空，便于问题排查。

\section{数据结构设计}

\subsection{持久化存储设计}
软件采用同步键值存储实现多端持久化（App、H5、小程序通用）。聊天会话以数组形式
存储，每个会话包含编号、标题、创建与更新时间、概要、压缩进度、设置快照、情景历史、
记忆与消息列表；另有当前会话编号用于记录活动会话。

设置项统一以前缀 chabot\_setting\_ 存储；记忆、情景历史、表情列表、API 预设、
背景图等按各自键名独立存储。首次启动会把旧版聊天历史自动迁移为第一个会话。

\subsection{主要数据结构}
\begin{itemize}[itemsep=0pt,parsep=0pt]
  \item \textbf{会话对象：}包含 id、title、created\_at、updated\_at、summary、
        compressedUntil、settings、scenes、memories、messages 字段；
  \item \textbf{消息对象：}包含 id、role、content、time 字段；
  \item \textbf{记忆对象：}包含 id、category、content、importance、level、
        created\_at、last\_accessed\_at、access\_count 字段；
  \item \textbf{表情对象：}包含 id、name、src、created\_at 字段；
  \item \textbf{情景对象：}包含 text、time 字段。
\end{itemize}
各对象字段随模块实现相应扩展。

\section{接口设计}

\subsection{外部接口}
\textbf{LLM 对话接口：}调用 OpenAI 兼容接口 /chat/completions，请求体包含 model、
messages 与 temperature、stream 等参数，强制非流式并解析响应 content 字段；
通过 reasoning\_effort 控制思考型模型的思考模式；服务端不识别参数时自动移除参数
降级重试。

\textbf{语音合成接口：}对接 Qwen-TTS 非流式语音合成接口，请求含模型名、音色与文本，
响应解析音频直链用于播放。

\subsection{内部接口}
软件将各域对外开放的公共接口通过存储门面与聊天门面统一导出，页面层仅依赖门面接口
即可完成全部业务，降低耦合。例如记忆增删改查、会话管理、设置读写、表情管理与解析、
日志记录等均通过门面暴露，调用方无需感知底层存储实现。

\section{技术特点}

\subsection{跨端一次编写、多端运行}
基于 uni-app（Vue 3）与条件编译，一套源代码即可打包为 App（Android/iOS）、H5 与
微信小程序，显著降低多平台开发与维护成本；针对各平台差异（数据持久化、图片压缩、
语音播放等）以条件编译与运行时能力检测适配。

\subsection{纯端侧、零后端}
聊天、记忆与设置全部在端上完成，不需要专属服务器；仅按需对接用户自配的大模型与
语音接口，最大限度保护用户数据隐私。

\subsection{三级记忆与多样化检索}
采用 L1/L2/L3 三级半衰期记忆模型、重要性与召回强化因子、MMR 多样性与新鲜槽混合
检索，实现个性化、可持续的记忆与对话。

\subsection{情景感知与上下文压缩}
借助 Scene 行感知用户当前情景，支持实时与虚拟时间模式；将历史对话压缩为概要注入
请求，兼顾上下文连贯与成本控制。

\subsection{拟真主动聊天}
通过链式定时器与持久化到期时刻实现低后台占用的主动消息调度，配合后台通知模拟真人
互动，并提供完整门禁与日志可观测机制。

\subsection{丰富的辅助能力}
内置自定义表情包、语音朗读、多会话、API 预设、调试日志面板、聊天记录导出等功能，
形成完整、易用的端侧聊天工具链。

\section{测试结果}

\subsection{测试环境}
测试在开发机（Windows 10）以 Node.js 运行自动化断言脚本，并在 Android 真机、
微信开发者工具、H5 浏览器进行手工冒烟测试。

\subsection{测试内容与结论}
针对核心逻辑编写断言测试：记忆系统（三级存储、时间衰减、去重合并、召回冷却、
分层检索、维护升降级、Memory 行解析）与表情包系统（名称校验、增删改、解析拆分、
占位提取、拖拽重排）均通过全部断言。

对聊天发送、回复格式校验、情景更新、记忆入库、上下文压缩、拟真调度、语音合成等
主流链路进行手工验证，功能符合设计预期，回复格式校验与记忆写入保持稳定。

跨端验证显示：App、H5、微信小程序三端核心功能一致，存储与运行正常，符合设计要求。

\section{使用说明}

\subsection{安装与导入}
开发者使用 HBuilderX 导入本项目工程，在开发环境配置大模型接口后，将工程运行到
浏览器、真机或模拟器；也可按平台打包发布。用户在设置页填写接口地址、API Key 与
模型名并保存，即可开始对话。

\subsection{功能使用}
\textbf{聊天：}在输入框输入消息发送；点击头部人格名可配置当前对话人格；历史弹窗
可切换、删除会话或压缩上文、导出对话；聊天记录超过设定条数时出现侧边滑块定位消息，
并支持一键回到底部与加载更早消息。

\textbf{记忆管理：}记忆页可查看、新建、编辑记忆内容、优先级与级别，支持筛选与
多选删除；记忆由 AI 自动维护，也支持手动管理。

\textbf{表情包：}表情管理页可批量上传并逐张命名、改名、删除与拖拽排序；
聊天中点击表情栏可插入 \$表情名\$ 占位。

\textbf{拟真聊天：}在设置中开启拟真聊天并设定活动时段与频率，AI 会主动发消息；
支持后台通知与调试发送。

\textbf{语音阅读：}在设置页开启语音阅读并填写语音接口 Key，AI 回复将以语音朗读，
表情不朗读。

\textbf{设置与调试：}设置页集中管理接口、思考模式、数据与背景等；调试日志面板
可查看请求、响应与错误详情，便于排查问题。

\section{界面与交互说明}

\subsection{聊天主界面}
聊天页顶部展示 AI 伙伴名称（点击可配置人格）、历史、新对话与清空按钮。中部为
消息列表，支持滚动、侧边滑块定位、一键回到底部，以及上翻加载更早消息；背景图片
支持自定义设置。底部为输入栏，右侧可展开表情栏插入表情，支持唤起键盘。情景条
随对话动态展示 AI 判断的用户当前情景，点击可编辑或清除。

\subsection{记忆管理界面}
记忆页以列表展示全部记忆，支持按关键字筛选、新建记忆、编辑内容与优先级级别、
开启多选模式批量删除，帮助用户直观掌控 AI 的个性化记忆。

\subsection{表情管理界面}
表情管理页支持批量上传图片并逐张命名，可改名、删除，支持长按拖动排序，实时同步
到聊天表情栏。

\subsection{设置界面}
设置页以分区形式组织：接口配置、思考模式、请求次数、人格配置、聊天表情包、
语音阅读、聊天背景、上下文压缩、拟真聊天、数据管理与调试日志等区块，并提供若干
设置子页（外观、接口、人格、拟真、语音、压缩、数据、日志、帮助、关于）。

\subsection{交互反馈}
界面统一使用主色 \#5b7cfa 与 rpx 自适应单位，弹窗共用统一的遮罩与面板样式；
列表数据复制副本渲染避免污染存储；消息列表渲染采用窗口化与稳定标识键，
保障表情较多场景下的流畅体验。

\section{异常处理与容错}

\subsection{接口容错}
LLM 请求失败或返回格式异常时，软件依据格式校验结果自动重新请求，最多重试设置项
指定的次数，达上限抛出错误并由页面提示；服务端不识别思考模式参数返回错误时自动
移除该参数降级重试。回复校验不通过时不写入历史，避免脏数据进入聊天与记忆。
语音首次播放失败会重新合成一次以获得新音频链接重试，兜住瞬时网络或解码问题。

\subsection{存储容错}
所有持久化写入均带异常兜底，避免写入失败抛出影响界面；存储层对平台差异
（App 与小程序存文件、H5 存 base64）做出适配与降级，兼容旧数据结构并自动迁移。

\subsection{调度容错}
拟真聊天调度使用链式定时器与持久化到期时刻，进程重启后自动对齐；每次到期无论
发送成败均记录日志（成功、跳过附原因、异常附堆栈），杜绝静默失败。

\subsection{数据安全}
聊天记录、记忆与设置全部保存在用户本地设备，不上传云端；导出对话记录与调试日志
均不包含接口密钥等敏感信息。

\section{版本记录}
\begin{itemize}[itemsep=0pt,parsep=0pt]
  \item \textbf{V1.0.0（初始版）：}完成聊天主链路、人格预设、三级记忆雏形、多会话、
        设置与数据持久化；
  \item \textbf{V1.1.x：}新增上下文压缩、情景感知、表情包系统与聊天背景图；
  \item \textbf{V1.2.x：}新增语音阅读（TTS）、上下文压缩优化、接口预设与思考模式控制；
  \item \textbf{V1.3.x：}新增拟真聊天与后台通知、回复格式严格校验、调试日志面板；
  \item \textbf{V1.4.x：}优化记忆系统（MMR 检索、容量管理）、消息列表渲染性能、
        导出与数据管理。
\end{itemize}
当前登记版本为 V1.4.3。各版本兼容演进，持续完善功能与稳定性。

\section{附录：主要程序文件说明}
\sloppy

\subsection{聊天服务层}
chat.js：聊天门面与发送主链路；chat-state.js：共享状态；chat-settings.js：
设置域；chat-conversations.js：会话域；chat-compress.js：压缩域；
chat-proactive.js：拟真聊天域。

\subsection{存储层}
storage.js：持久化门面；storage-state.js：存储底层与共享状态；
storage-conversations.js：会话域；storage-scene.js：情景域；storage-api.js：
API 预设域；storage-background.js：背景图域。

\subsection{记忆与提示词}
memory.js：记忆核心；memory-similarity.js：相似度；memory-parse.js：Memory 行解析；
memory-constants.js：常量；prompts.js：系统提示词与人格预设。

\subsection{辅助能力}
emojis.js：表情数据层；llm.js：LLM 客户端；tts.js：语音阅读；log.js：调试日志；
notify.js：系统通知；export.js：记录导出。

\subsection{页面层}
pages/chat/chat.vue 及 components 目录：聊天页与七个子组件；
pages/memory/memory.vue：记忆页；pages/emoji/emoji.vue：表情页；
pages/settings 目录：设置页与各设置子页。

\subsection{测试脚本}
scripts/test-memory.mjs：记忆核心逻辑断言测试，覆盖三级存储、时间衰减、去重合并、
召回冷却、分层检索与维护升降级等场景；scripts/test-emojis.mjs：表情包逻辑断言测试，
覆盖名称校验、增删改、解析拆分、占位提取与重排等场景。两个脚本均在 Node 环境
通过 npm test 运行，对核心纯函数逻辑进行自动化回归验证。

\subsection{工程配置与入口文件}
pages.json：页面注册、窗口样式与 tabBar 配置；manifest.json：应用标识、版本号、
平台发布信息与权限声明；App.vue：应用入口，负责启动初始化、存储迁移与前后台
状态同步；main.js：框架入口，创建应用实例并挂载。

\section{结语}
ChabotAI伙伴 V1.4.3 以“无后端、纯端侧、跨平台”为核心设计理念，围绕大语言模型
构建了集智能对话、三级记忆、情景感知、上下文压缩、表情包互动、语音阅读与拟真
主动聊天于一体的完整端侧聊天应用。软件在架构上采用门面与域模块分层设计，
在工程上以自动化断言测试保障核心逻辑质量，在体验上兼顾多端一致性与平台特性，
形成了良好的可扩展性与可维护性。本文档与随附源程序共同构成软件的完整技术资料，
供登记审查与后续版本演进参考。

\vspace{1em}
\noindent
\begin{tabular}{@{}ll@{}}
  软件名称： & ChabotAI伙伴 \\
  软件版本： & V1.4.3 \\
  著作权人： & 陈聪 \\
  文档类型： & 软件设计说明书 \\
  编制日期： & 2026 年 9 月 \\
\end{tabular}

\end{document}